\chapter{Conclusões e Trabalhos Futuros}
\label{chap:conclusoes-e-trabalhos-futuros}

Este trabalho construiu, validou e comparou uma progressão incremental de seis formulações de fluxo de potência AC, desenvolvidas em \textit{JuMP/Julia} sobre o \textit{solver} Ipopt. Duas formulações foram tomadas como referência: o FPO clássico (\texttt{solve\_ac\_opf}) e o FP clássico (\texttt{solve\_ac\_pf}) do \textit{PowerModels.jl}. As outras quatro foram construídas \textit{do zero} para incorporar, sequencialmente, as ações de controle do programa ANAREDE: limites de geração reativa (QLIM) e de magnitude de tensão (VLIM) na formulação F2; controle de bancos \textit{shunt} chaveáveis (CSCA) na F3; controle automático de \textit{tap} de transformadores (CTAP) na F4; e um esquema hierárquico de corte/alívio de carga (rotulado como DERA, em referência ao nome do \textit{script}) na F5. Em todas as formulações \textit{hand-built}, os controles foram introduzidos por meio de variáveis de folga penalizadas em uma função objetivo do tipo \textit{soft-constraint}, evitando a representação descontínua dos laços de \textit{type-switching} do ANAREDE.

\section{Síntese das contribuições}
\label{sec:contribuicoes-do-trabalho}

A bateria padronizada de vinte e sete casos de teste, que cobre redes pedagógicas de três a nove barras, redes de médio porte (300 e 500 barras), recortes reduzidos do Nordeste do SIN (\texttt{caso\_red} e \texttt{caso\_red2}) e o \textit{snapshot} integral do SIN com 13\,338 barras (\texttt{CASO\_VER\_MAXDIU.PWF} do PAR/PEL 2027-2031), forneceu evidência empírica de três contribuições principais.

A primeira é que as formulações \textit{hand-built} F2 a F5 convergem em uma faixa estritamente maior de casos do que as referências \textit{PowerModels}: vinte e três casos contra dezoito da F1 e quinze da F0. Em particular, em oito casos (\texttt{4busfrank\_vlim}, \texttt{5busfrank\_csca}, \texttt{test\_system}, \texttt{caso\_red2}, \texttt{300bus}, \texttt{3bus\_DCline}, \texttt{3bus\_DSHL} e \texttt{caso\_red}) ao menos uma das formulações \textit{hand-built} encontrou uma solução factível, enquanto ao menos uma das referências não convergiu. Esse comportamento, esperado pela construção \textit{soft-constraint}, valida a tese de que a representação dos controles ANAREDE via folgas penalizadas é uma alternativa numericamente robusta às heurísticas iterativas do programa de origem.

A segunda contribuição é o agrupamento sistemático das soluções em \textit{clusters} de equivalência numérica que reproduzem o efeito esperado de cada controle: o cluster $\{\text{F1,F2}\}$ aparece quando QLIM/VLIM saturam em zero; o cluster $\{\text{F3,F4}\}$ ou $\{\text{F3}\}\cdot\{\text{F4}\}$ aparece em redes onde CSCA ou CTAP é ativo; e o DERA (formulação F5) comporta-se como F4 na maioria dos casos pequenos (porque o \textit{solver} prefere aceitar pequenas folgas de tensão a cortar carga) e forma cluster próprio em redes com margem reativa estreita, em que o corte hierárquico passa a ser efetivamente acionado. Esse padrão fornece um \textit{fingerprint} computacional que pode ser usado, em trabalhos futuros, para detectar regressões em implementações controladas.

A terceira contribuição é negativa, mas relevante: nenhuma das seis formulações converge sobre o \textit{snapshot} integral do SIN (\texttt{CASO\_VER\_MAXDIU.PWF}) dentro do orçamento de \texttt{max\_iter = 3000} iterações e do \textit{wall time} permitido. As formulações F0, F2 e F4 atingem o limite de tempo (\textbf{KIL}) sem sinalizar infactibilidade, enquanto F1, F3 e F5 terminam em LOCALLY\_INFEASIBLE. A análise dos logs sugere que o ponto de partida \textit{flat} está suficientemente longe de uma solução factível para que o método de pontos interiores do Ipopt não saia da região de inviabilidade dentro do orçamento. Esse resultado delimita com clareza a fronteira atual da abordagem \textit{soft-constraint} pura sem inicialização guiada. Vale notar, contudo, que os dois recortes do mesmo caso (\texttt{caso\_red} com 2\,599 barras e \texttt{caso\_red2} com 1\,033 barras), descritos na \autoref{sec:casos-reduzidos}, convergem em F2 a F5 e ilustram em escala intermediária o ganho de robustez das formulações \textit{hand-built} sobre as referências.

\section{Limitações}
\label{sec:limitacoes}

As principais limitações deste trabalho são três. Primeiro, todas as formulações tratam $b^{sh}_s$ (CSCA) e $t_{ij}$ (CTAP) como variáveis contínuas, embora a operação real seja discreta. A discretização posterior por \textit{rounding} ou em formulação MINLP não foi realizada, e o impacto da relaxação contínua sobre a fidelidade ao ANAREDE não foi quantificado. Segundo, o esquema de corte de carga rotulado como DERA na formulação F5 é uma versão contínua e suave de um mecanismo originalmente discreto e em emergência (do tipo ERAC do ONS): o \textit{solver} pode cortar frações arbitrárias de carga em qualquer barra, sem respeitar blocos pré-definidos nem temporização. Terceiro, os elos HVDC são tratados com fluxos fixados em vez de variáveis livres. Essa é uma decisão correta sob a hipótese de \textit{fluxo base já dado}, mas que limita a aplicabilidade da formulação a cenários em que o despacho HVDC pode ser otimizado.

\section{Trabalhos futuros}
\label{sec:trabalhos-futuros}

A continuidade natural deste trabalho passa por três frentes complementares. A primeira é viabilizar a convergência sobre o caso SIN. Duas possibilidades concretas são: (i) inicialização \textit{warm-start} baseada na solução de uma formulação relaxada (DC ou SOC) sobre o mesmo caso, em substituição ao \textit{flat start}; e (ii) reformulação do problema com escalonamento explícito de penalidades $\rho$, em vez do valor único $10^6$, com aumento progressivo até a estabilização das folgas. Ambas as estratégias são bem documentadas em \cite{molzahn2019survey} e mantêm o \textit{solver} Ipopt como motor. A primeira possibilidade foi parcialmente validada já neste trabalho pela formulação auxiliar FDC (\autoref{subsec:fdc-resultado}), que aplicou \texttt{solve\_dc\_pf} do \textit{PowerModels} sobre o mesmo \texttt{CASO\_VER\_MAXDIU.PWF} e convergiu em poucas iterações com status \texttt{LOCALLY\_SOLVED}, confirmando que o obstáculo das formulações AC sobre o SIN é numérico e não estrutural. Resta, como passo natural, encadear o resultado DC como \texttt{start} das variáveis $V$ e $\theta$ das formulações F2 a F5, com prévia regularização dos ângulos (que na solução DC chegam a $\pm 170^\circ$, fora do regime de validade da linearização).

A segunda é a discretização dos controles CSCA e CTAP, transformando as formulações F3 e F4 em problemas MINLP. O ecossistema Julia já oferece \textit{solvers} adequados (Juniper, Alpine, SCIP) acoplados ao JuMP, e o \textit{benchmark} sobre a mesma bateria permitiria quantificar o gap entre a relaxação contínua e a operação discreta real.

A terceira é o acoplamento direto das formulações F2 a F5 ao pacote \textit{ControlPowerFlow.jl}, em desenvolvimento no grupo de pesquisa do orientador. As formulações \textit{hand-built} apresentadas aqui podem servir como \textit{oracle} de validação para a implementação modular do pacote, ao mesmo tempo em que o pacote pode oferecer uma API mais limpa para os controles agora codificados linha a linha em cada \textit{script}. O trabalho de \cite{chavarry2024} estabelece a referência metodológica sobre a qual essa integração pode ser construída.
